%! Comparing Distributions — Unit 2, Lesson 2.4 — Student Packet Cover Sheet
% THREE packet rows — Warm-Up, Guided Notes & Practice, Homework. The homework is the individual
% practice, assigned at Close & Assign and done outside class (Unit 2 on), and it is SCORED, so
% its cell is a \blank{}, never NA.
\documentclass[12pt]{article}
\usepackage{saar-article}
\usepackage{saar-boxes}
\usepackage{ltablex}
\keepXColumns
% Measured banner: the band grows to fit the title block, so a title long enough to wrap grows
% the band rather than running out of it (the stats course's \coverbanner; saar-article does not
% carry one, so it is defined here — every cover copies this block verbatim).
\newsavebox{\bannerbox}
\newlength{\bannerht}
\newlength{\coverbannerpad}
\setlength{\coverbannerpad}{0.16in}       % breathing room above and below the text
\newcommand{\coverbanner}[2]{%
  \sbox{\bannerbox}{%
    \begin{minipage}{\dimexpr\paperwidth-1.4in\relax}%
      \centering
      {\color{white}\textbf{\LARGE \CourseName}}\\[4pt]
      {\color{frostmid}\large\textbf{#1}}\\[6pt]
      {\color{white}\textbf{\large #2}}%
    \end{minipage}}%
  \setlength{\bannerht}{\dimexpr\ht\bannerbox+\dp\bannerbox+2\coverbannerpad\relax}%
  \noindent\begin{tikzpicture}[remember picture, overlay]
    \fill[cerulean] (current page.north west)
      rectangle ([yshift=-\bannerht]current page.north east);
    \node[anchor=north, inner sep=0pt]
      at ([yshift=-\coverbannerpad]current page.north) {\usebox{\bannerbox}};
  \end{tikzpicture}\par
  % The text body starts 0.6in below the page edge (geometry top=0.6in), so reserve only what
  % the band overruns that, plus a gap under the band.
  \vspace*{\dimexpr\bannerht-0.6in+0.16in\relax}%
}

\begin{document}

% Full-bleed cerulean banner, sized to the title block.
\coverbanner{Unit 2}{Lesson 2.4 \quad Comparing Distributions}

\vspace{0.06in}
% NAMESTRIP: this is the ONLY component that carries the name row.
\namedateperiod
\vspace{0.18in}

% GRADUAL RELEASE: the vocabulary is taught in the guided notes, so the targets name it
% plainly. The standard codes (PS.DS.3a-b) stay in the teacher-facing plan.
\begin{learningtargetbox}
\small
\begin{enumerate}[leftmargin=1.5em, itemsep=5pt]
  \item \ldots{} read a \textbf{back-to-back stemplot}, \textbf{parallel dot plots}, and
        \textbf{parallel boxplots} --- and compare groups only on a \textbf{common scale}.
  \item \ldots{} write a \textbf{comparison} of two distributions in context, using
        comparative words for shape, outliers, center, and spread.
  \item \ldots{} choose the measures that fit the shape: median and IQR when a distribution is
        skewed or has an outlier, mean and standard deviation when both are roughly symmetric.
  \item \ldots{} decide which group is more \textbf{consistent} when the centers are the same
        --- and explain why the spread, not the center, decides.
\end{enumerate}
\end{learningtargetbox}

\vspace{0.14in}

\begin{tocbox}
\small
\renewcommand{\arraystretch}{1.5}
\begin{tabularx}{\linewidth}{@{} c l X r @{}}
\rowcolor{frost}
\textbf{\#} & \textbf{Component} & \textbf{Description} & \textbf{Score} \\
\midrule
1 & Warm-Up & Tony's Pizza delivery times: quartiles, IQR, and the $1.5 \cdot \text{IQR}$
             fence --- then two tiny lists with the same mean & \blank{1.2cm} \\
2 & Guided Notes \& Practice & Back-to-back stemplots and parallel displays on a common
             scale; comparing shape, outliers, center, and spread in context; two pizza shops
             with the same median --- then paper airplanes we work together & \blank{1.2cm} \\
% Row 3 is the lesson's GRADED practice, assigned at the close. Packet pages are the default; a
% Desmos activity may replace them on a given day, decided at Close & Assign. The row and its
% score cell never change.
3 & Homework & Volt vs.\ Nova phone batteries: five-number summaries, the right measures, and
             which brand for which buyer --- \textbf{scored, assigned today, due the first
             class after two study halls} & \blank{1.2cm} \\
\midrule
  & \multicolumn{2}{r}{\textbf{Total}} & \blank{1.2cm} \\
\end{tabularx}
\end{tocbox}

\vspace{0.14in}

% Keep in Mind carries the lesson's CONTENT — the rule, the test, the trap — never its process.
\begin{remindbox}
\small
A \textbf{comparison} of two distributions compares all four features --- shape, outliers,
center, and spread --- with comparative words (\emph{greater than, less than, about the same
as}), in context, and only on a \textbf{common scale}. Use median and IQR when either group is
skewed or has an outlier; mean and SD when both are roughly symmetric. \textbf{Watch out:} two
groups with the same center are not the same. When the centers match, the spread decides ---
the group with less spread is more \textbf{consistent}. Two separate descriptions side by side
are not a comparison.
\end{remindbox}

\end{document}
